\documentclass[a4paper, 11pt]{article}
\usepackage[spanish, es-tabla]{babel}
\usepackage[utf8]{inputenc}
\usepackage{array}
\usepackage{booktabs}      % Para \toprule, \midrule, \bottomrule
\usepackage{float}         % Para usar [H] en tablas y figuras
\usepackage{mdwmath}
\usepackage{mathabx}
\usepackage{subcaption}
\usepackage{siunitx}
\usepackage{cite}
\usepackage{caption}
\usepackage{algorithm}
\usepackage[pdftex]{graphicx}
\usepackage{amsmath}
\usepackage[table]{xcolor}
\usepackage{enumitem}
\usepackage{listings}
\usepackage{xcolor}

\definecolor{codegreen}{rgb}{0,0.6,0}
\definecolor{codegray}{rgb}{0.5,0.5,0.5}
\definecolor{codepurple}{rgb}{0.58,0,0.82}
\definecolor{codeblue}{rgb}{0.0,0.2,0.8}

\lstdefinestyle{estiloArduinoCompacto}{
    language=C++,
    backgroundcolor=\color{white},   
    commentstyle=\color{codegreen},
    keywordstyle=\color{codeblue},
    numberstyle=\tiny\color{codegray},
    stringstyle=\color{codepurple},
    basicstyle=\ttfamily\tiny, 
    breakatwhitespace=false,         
    breaklines=true,                 
    captionpos=b,                    
    keepspaces=true,                 
    numbers=left,                    
    numbersep=3pt,              
    showspaces=false,                
    showstringspaces=false,
    showtabs=false,                  
    tabsize=1,
    frame=tb, 
    xleftmargin=1.5em,
    framexleftmargin=1.5em
}

\lstset{style=estiloArduinoCompacto}

\usepackage{multicol}
\usepackage[breaklinks]{hyperref}
\hypersetup{
    colorlinks=true,
    linkcolor=black,
    citecolor=blue,
    urlcolor=blue
}
\usepackage{xurl}

\usepackage[left=2cm, right=2cm, top=2cm, bottom=2cm]{geometry}

\setlength{\parindent}{0pt}

\usepackage{titlesec}

% --- CORRECCIÓN PARA QUE REFERENCIAS SEA UNA SECCIÓN NUMERADA ---
\usepackage{etoolbox}

\makeatletter
\patchcmd{\thebibliography}
  {\section*{\refname}}
  {\section{\refname}}
  {}{}
\makeatother

\begin{document}

\begin{center}
    \huge{\textbf{INFORME N°2: Aplicación De Machine Learning En La Agricultura Alimentada por paneles solares}}\\
    \vspace{3mm} 
    \large \textbf{IF- Ingeniería Solar} \\
    \vspace{1mm}
    Olivas Huamán Dante, dante.olivas.h@uni.pe \\
              
    \vspace{1mm}
    Universidad Nacional de Ingeniería, Facultad de Ciencias \\
    \vspace{1mm}
    19/05/2026
\end{center}

\vspace{2mm}

\hrule height 1.2pt 
\vspace{1mm}

\begin{abstract}
Este informe detalla el diseño, implementación y dimensionamiento energético del Proyecto Aripura (AgroPredict), una plataforma tecnológica de agricultura de precisión orientada a zonas remotas y comunidades altoandinas de escasos recursos en el Perú. El sistema consta de un nodo sensor basado en el microcontrolador de bajo costo ESP32 que recopila variables ambientales (temperatura, humedad relativa, presión barométrica, humedad del suelo y luminosidad) y las transmite localmente mediante HTTP/JSON a un servidor perimetral. En este servidor se ejecuta un modelo de Machine Learning basado en el algoritmo \textit{HistGradientBoostingClassifier} que realiza inferencias en tiempo real para diagnosticar estados críticos del cultivo, tales como estrés hídrico severo o leve, asfixia radicular, alertas fúngicas y de plagas, así como estrés térmico por heladas inminentes. Todo el ecosistema técnico, incluyendo el nodo de sensores, la computadora local de procesamiento y la antena satelital Starlink que asegura la conectividad en áreas sin infraestructura de telecomunicaciones tradicionales, se alimenta a través de un sistema solar fotovoltaico aislado (off-grid) que adicionalmente cubre las necesidades eléctricas domésticas del agricultor. Se presenta el dimensionamiento energético del sistema fotovoltaico con un requerimiento diario de \SI{1962}{Wh/dia}, demostrando la viabilidad técnica, el valor agregado para la prevención de pérdidas agrícolas y el impacto socioeconómico en la resiliencia del sector rural peruano.
\end{abstract}
\vspace{2mm}

\hrule height 1.2pt
\vspace{5mm}

\begin{multicols}{2}

\section{Introducción}
La geografía del Perú presenta una notable diversidad climática, pero también aísla a miles de comunidades agrícolas en zonas altoandinas y amazónicas. En estas áreas remotas, los agricultores se enfrentan de forma constante a los efectos del cambio climático, manifestados en heladas súbitas, sequías cíclicas e infestación atípica de plagas \cite{ref1}. La falta de acceso a información meteorológica local y de precisión, sumada a la ausencia de redes de conectividad tradicionales (telefonía móvil y fibra óptica), restringe severamente la capacidad de los productores para anticipar estos desastres y optimizar sus recursos hídricos y nutricionales.

En este contexto surge el \textbf{Proyecto Aripura} (bajo el desarrollo del sistema de software \textit{AgroPredict}), concebido como una solución integral que introduce el Internet de las Cosas (IoT) y la inteligencia artificial directamente en el campo rural. El sistema no depende de infraestructura externa de energía o comunicaciones: se autoabastece mediante paneles solares fotovoltaicos y aprovecha la órbita terrestre baja por medio de la tecnología Starlink para proveer internet de banda ancha de manera ininterrumpida. Esto permite habilitar un canal de datos bidireccional mediante el cual el agricultor puede recibir alertas automáticas, descargar datos analíticos y compartir información con cooperativas y especialistas agrónomos a nivel nacional.

La presente investigación detalla tanto el marco teórico como la implementación práctica de la telemetría IoT del Proyecto Aripura. Asimismo, se expone el motor de \textit{Machine Learning} (ML) entrenado con variables sintéticas de alta fidelidad agronómica, y se realiza el dimensionamiento formal del sistema de energía solar necesario para mantener operando de manera continua la antena Starlink, el nodo sensor y las necesidades domésticas de la vivienda del agricultor.

\section{Objetivos}

\subsection{Objetivo general}
Diseñar, dimensionar y evaluar un sistema de monitoreo agrícola de precisión autoalimentado por energía solar fotovoltaica para zonas rurales del Perú, el cual implementa modelos de Machine Learning sobre variables microclimáticas y conectividad satelital Starlink para mitigar riesgos en los cultivos y mejorar la toma de decisiones.

\subsection{Objetivos específicos}
\begin{itemize}
    \item Implementar un nodo de adquisición de datos microclimáticos basado en el microcontrolador ESP32 utilizando sensores de temperatura, humedad del aire, presión atmosférica, humedad capacitiva del suelo e irradiancia relativa.
    \item Desarrollar y evaluar un modelo de clasificación multiclase basado en el algoritmo \textit{HistGradientBoosting} para identificar estados patológicos y fisiológicos del cultivo.
    \item Dimensionar el sistema fotovoltaico aislado y el banco de almacenamiento de baterías de litio ($LiFePO_4$) que soporte la carga del nodo, la antena Starlink y el consumo doméstico básico.
    \item Estructurar el motor de ingeniería de características (\textit{Feature Engineering}) para el cálculo del Déficit de Presión de Vapor (VPD) y la acumulación de Grados-Día (GDD) en tiempo real.
\end{itemize}

\section{Marco teórico}

\subsection{Adquisición Microclimática y Sensores IoT}
La precisión de los modelos predictivos depende de la calidad física de los datos de entrada. En el Proyecto Aripura se integran cuatro sensores clave:
\begin{enumerate}
    \item \textbf{DHT11 (Humedad y Temperatura del Aire):} Utiliza un componente resistivo para la humedad y un termistor NTC para la temperatura. Proporciona una precisión de $\pm 2^\circ\text{C}$ y $\pm 5\%$ HR, ideal para propósitos de validación general.
    \item \textbf{BMP180 (Presión Barométrica y Clima de Precisión):} Transductor piezoeléctrico de alta precisión con comunicación $I^2C$ que permite lecturas exactas de la presión atmosférica ($Pa$) y temperatura de respaldo, fundamental para modelar la altitud y los gradientes térmicos locales \cite{ref2}.
    \item \textbf{Sensor de Humedad del Suelo Capacitivo:} A diferencia de los sensores resistivos convencionales, mide la permitividad dieléctrica del suelo mediante un campo eléctrico alterno. Esto evita la corrosión galvánica de los electrodos en contacto con el suelo húmedo y garantiza lecturas estables a largo plazo \cite{ref3}. El valor se obtiene como lectura analógica de 12 bits (\num{0}-\num{4095}).
    \item \textbf{Fotorresistencia (LDR):} Su resistencia disminuye conforme aumenta la radiación solar incidente, permitiendo estimar la irradiancia relativa a través de un divisor de voltaje conectado a un pin ADC.
\end{enumerate}

\subsection{Fisiología Vegetal e Indicadores Agronómicos}
El modelo de Machine Learning no solo recibe datos crudos, sino también variables sintéticas que modelan directamente las interacciones planta-atmósfera:

\textbf{Déficit de Presión de Vapor (VPD):} Es la diferencia entre la cantidad de humedad que el aire puede retener cuando está saturado y la cantidad de humedad real en ese momento. Se calcula en kilopascales ($kPa$) a través de la ecuación de Tetens-Buck:
\begin{equation}
    SVP = 0.6108 \cdot \exp\left(\frac{17.27 \cdot T}{T + 237.3}\right)
\end{equation}
\begin{equation}
    AVP = SVP \cdot \left(\frac{RH}{100}\right)
\end{equation}
\begin{equation}
    VPD = SVP - AVP
\end{equation}
Donde $T$ es la temperatura en $^\circ\text{C}$, $RH$ es la humedad relativa en $\%$, $SVP$ es la presión de vapor de saturación y $AVP$ es la presión de vapor real. Un VPD bajo ($<0.4\,kPa$) indica que el aire está saturado, lo que inhibe la transpiración de la planta y favorece la proliferación fúngica; un VPD alto ($>1.5\,kPa$) induce el cierre estomático para evitar la deshidratación, desencadenando estrés hídrico \cite{ref4}.

\textbf{Punto de Rocío ($T_{dp}$):} Temperatura a la cual el vapor de agua en el aire comienza a condensarse en forma líquida sobre las hojas. Se calcula mediante la fórmula de Magnus:
\begin{equation}
    \alpha(T, RH) = \frac{17.27 \cdot T}{237.3 + T} + \ln\left(\frac{RH}{100}\right)
\end{equation}
\begin{equation}
    T_{dp} = \frac{237.3 \cdot \alpha(T, RH)}{17.27 - \alpha(T, RH)}
\end{equation}
La presencia de agua líquida libre combinada con temperaturas moderadas ($15^\circ\text{C}$ - $25^\circ\text{C}$) activa esporas de patógenos como \textit{Botrytis cinerea} o \textit{Mildiu}.

\textbf{Grados-Día Acumulados (GDD):} Es una medida de la acumulación térmica utilizada para predecir el desarrollo fisiológico de cultivos y plagas (fenología). Se define como:
\begin{equation}
    GDD = \sum \max(T - T_{base}, 0) \cdot \Delta t
\end{equation}
Donde $T_{base}$ es la temperatura mínima a partir de la cual el organismo detiene su desarrollo (para este MVP se fijó $T_{base} = 10^\circ\text{C}$), y $\Delta t$ es la porción del día en horas transcurridas.

\subsection{Machine Learning: Gradient Boosting en el Edge}
Para clasificar los estados biológicos del cultivo se emplea el algoritmo \textit{HistGradientBoostingClassifier}. Este algoritmo de scikit-learn, inspirado en LightGBM, discretiza las características continuas en contenedores enteros (\textit{bins}) de tamaño fijo. Esto reduce la complejidad computacional del cálculo de divisiones óptimas de $O(N \log N)$ a $O(N)$, permitiendo un entrenamiento y una inferencia sumamente rápidos que pueden correr en hardware perimetral de bajo costo, tales como mini-PCs o laptops reutilizadas presentes en la vivienda del agricultor.

\subsection{Sistemas Solares Fotovoltaicos y Conectividad Starlink}
La operación continua en zonas aisladas requiere el uso de energía solar. Un sistema fotovoltaico aislado clásico consta de paneles fotovoltaicos, un controlador de carga de tecnología MPPT (\textit{Maximum Power Point Tracking}), un banco de almacenamiento de baterías (preferentemente de química de Litio-Hierro-Fosfato, $LiFePO_4$, por su alta densidad energética y vida útil de más de 3000 ciclos), y un inversor de onda senoidal pura para alimentar la antena Starlink y electrodomésticos en corriente alterna ($220\,\text{V}$).

La antena Starlink V3 (de motor plano y alto rendimiento) consume de forma continua entre \SI{75}{W} y \SI{100}{W} debido al procesamiento de señales satelitales en banda Ka/Ku y el sistema de calefacción automática contra heladas/nieve. Este dispositivo representa la carga crítica dominante del sistema fotovoltaico.

\section{Metodología}

\subsection{Equipo utilizado}
El hardware seleccionado y validado físicamente para el Proyecto Aripura comprende:
\begin{itemize}
    \item \textbf{Unidad de Control:} Microcontrolador ESP32-WROOM-32E (doble núcleo, WiFi y Bluetooth integrados).
    \item \textbf{Sensor de Humedad Suelo:} Sensor capacitivo analógico v1.2.
    \item \textbf{Sensor Climatológico 1:} DHT11 de tres pines para humedad relativa.
    \item \textbf{Sensor Climatológico 2:} BMP180 de Bosch sobre bus I2C para presión barométrica.
    \item \textbf{Sensor de Iluminación:} Fotorresistencia LDR en configuración de divisor de tensión.
    \item \textbf{Servidor de Procesamiento Local:} Laptop de bajo consumo de energía (Intel Celeron, RAM 8GB) configurada con FastAPI bajo Python 3.10.
    \item \textbf{Módulo de Comunicación Satelital:} Kit de internet satelital Starlink Gen 3.
    \item \textbf{Subsistema Solar Fotovoltaico:} Panel solar monocristalino de \SI{500}{Wp}, controlador MPPT 30A, batería $LiFePO_4$ de \SI{24}{V} y \SI{120}{Ah}, e inversor de \SI{1000}{W} a \SI{220}{VAC}.
\end{itemize}

\subsection{Procedimiento experimental}
El flujo lógico del sistema se divide en tres fases principales: adquisición de datos, transmisión inalámbrica local y procesamiento perimetral basado en inteligencia artificial.

El microcontrolador ESP32 se configura para operar en modo \textit{Access Point} (AP), creando una red local segura llamada \texttt{Agro\_Kipura\_Tech}. Esto permite que el servidor local de bajo consumo (instalado en la vivienda del agricultor) se conecte de forma directa al microcontrolador. Al ESP32 se le asigna la dirección IP estática \texttt{192.168.4.1}, mientras que al servidor perimetral se le asigna mediante configuración de red la IP fija \texttt{192.168.4.2}.

Cada 5 segundos, el ESP32 despierta de su ciclo, realiza lecturas de los puertos analógicos (humedad del suelo en pin 34 y LDR en pin 32) e interactúa mediante la biblioteca \texttt{Wire} con el sensor BMP180 y mediante el protocolo digital de un solo hilo con el DHT11. A continuación, el ESP32 encapsula la información en un objeto JSON y la transmite al servidor local a través de una petición HTTP POST en el endpoint \texttt{/api/v1/telemetry}. En el Listado \ref{lst:codigoESP32} se presenta la sección central de adquisición y transmisión del código en C++.

\begin{lstlisting}[language=C++, caption={Función de envío de datos del ESP32 mediante HTTP POST.}, label={lst:codigoESP32}]
void sendDataToServer(float temp, float hum, float pres, float soil_analog, float light_analog) {
    HTTPClient http;
    http.begin(api_url);
    http.addHeader("Content-Type", "application/json");

    StaticJsonDocument<256> doc;
    doc["device_id"] = device_id;
    doc["temperature_c"] = temp;
    doc["humidity_air_pct"] = hum;
    doc["pressure_pa"] = pres;
    doc["soil_m_analog"] = soil_analog;
    doc["light_analog"] = light_analog;

    String requestBody;
    serializeJson(doc, requestBody);

    Serial.println("-> Enviando datos a Backend...");
    int httpResponseCode = http.POST(requestBody);

    if(httpResponseCode > 0) {
        String response = http.getString();
        Serial.print("<- Respuesta Servidor: ");
        Serial.println(response); 
    } else {
        Serial.print("[ERROR] Codigo: ");
        Serial.println(httpResponseCode);
    }
    http.end();
}
\end{lstlisting}

Una vez recibida la telemetría en la API del backend construida con FastAPI, se activa el motor de ingeniería de características. Se calculan dinámicamente las variables agronómicas $VPD$ (Eq. 3) y $T_{dp}$ (Eq. 5). Adicionalmente, el sistema mantiene en caché de memoria el historial temporal reciente para calcular variables acumulativas: los Grados-Día acumulados (Eq. 6) y las horas continuas con humedad relativa superior al $85\%$ y VPD inferior a $0.4\,kPa$.

Finalmente, estas 8 características estructuradas se introducen en el modelo cargado \texttt{agropredict\_model.pkl}. El modelo devuelve una predicción del estado de salud del cultivo, la cual es almacenada en un archivo de persistencia CSV (\texttt{telemetry\_log.csv}) de forma que pueda exportarse directamente a una hoja de cálculo en Excel para que el agricultor analice el progreso sin necesidad de conexión activa a internet. Si el internet satelital Starlink está encendido, el sistema puede sincronizar este archivo plano con bases de datos en la nube para permitir el monitoreo remoto por parte de agrónomos externos.

\section{Cálculos y Dimensionamiento}

\subsection{Datos experimentales y de carga}
Para el diseño y dimensionamiento del sistema solar fotovoltaico aislado que sustentará el Proyecto Aripura en una comunidad remota de la sierra del Perú (ej. Puno, Cusco o Cajamarca), se listan las cargas eléctricas de consumo continuo en la Tabla \ref{tab:cargas}.

\begin{table}[H]
\centering
\caption{Consumo energético diario estimado del Proyecto Aripura.}
\label{tab:cargas}
\begin{tabular}{lcccr}
\toprule
\textbf{Carga} & \textbf{Potencia} & \textbf{Uso} & \textbf{Voltaje} & \textbf{Energía} \\
& (W) & (h/día) & & (Wh/día) \\
\midrule
Antena Starlink & 75.0 & 24.0 & 220 VAC & 1800.0 \\
Servidor Local & 15.0 & 6.0 & 19 VDC & 90.0 \\
Nodo ESP32 & 0.5 & 24.0 & 5 VDC & 12.0 \\
Luz LED Doméstica & 10.0 & 5.0 & 220 VAC & 50.0 \\
Carga de Móvil & 10.0 & 1.0 & 5 VDC & 10.0 \\
\midrule
\textbf{Total diario} & & & & \textbf{1962.0} \\
\bottomrule
\end{tabular}
\end{table}

\subsection{Dimensionamiento Fotovoltaico}
El cálculo del dimensionamiento del generador fotovoltaico se rige por la siguiente ecuación:
\begin{equation}
    P_{fv} = \frac{E_d}{HSP \cdot \eta_{inst}}
\end{equation}
Donde:
\begin{itemize}
    \item $E_d$ es el requerimiento energético diario (\SI{1962}{Wh/dia}).
    \item $HSP$ representa las Horas Solar Pico promedio de la zona. Para la sierra peruana se asume un valor conservador de \SI{5.5}{h} (en invierno o épocas de lluvias).
    \item $\eta_{inst}$ es la eficiencia de instalación global, que considera pérdidas en el cableado, rendimiento del controlador de carga MPPT, polvo en el panel y calentamiento de celdas. Se adopta un estándar de $\eta_{inst} = 0.75$ ($75\%$).
\end{itemize}

Sustituyendo los valores de diseño:
\begin{equation*}
    P_{fv} = \frac{1962}{5.5 \cdot 0.75} = \frac{1962}{4.125} = 475.63 \text{ Wp}
\end{equation*}
Por consiguiente, se selecciona comercialmente \textbf{un (1) panel solar monocristalino de \SI{500}{Wp}} o, alternativamente, dos paneles de \SI{250}{Wp} conectados en serie para maximizar el rango de tensión de entrada del controlador MPPT.

\subsection{Dimensionamiento del Banco de Baterías}
La capacidad requerida del banco de almacenamiento debe asegurar la autonomía del sistema durante al menos un día nublado completo sin sol ($N = 1$ día de autonomía), protegiendo la vida útil de las baterías de litio $LiFePO_4$:
\begin{equation}
    Cap_{\text{bat}} = \frac{E_d \cdot N}{V_{\text{sistema}} \cdot DOD \cdot \eta_{\text{inv}}}
\end{equation}
Donde:
\begin{itemize}
    \item $V_{\text{sistema}}$ es el voltaje nominal del banco de baterías, fijado en \SI{24}{V} para reducir pérdidas por efecto Joule y sección de cable.
    \item $DOD$ es la profundidad de descarga permitida (para litio $LiFePO_4$ se establece de forma segura en $80\%$ o $0.80$).
    \item $\eta_{\text{inv}}$ es el rendimiento del inversor de corriente de \SI{24}{VDC} a \SI{220}{VAC}, estimado en $90\%$ o $0.90$.
\end{itemize}

Sustituyendo los parámetros:
\begin{equation*}
    Cap_{\text{bat}} = \frac{1962 \cdot 1}{24 \cdot 0.80 \cdot 0.90} = \frac{1962}{17.28} = 113.54 \text{ Ah}
\end{equation*}
Se opta por configurar un banco comercial de baterías de litio $LiFePO_4$ con una capacidad estándar de \textbf{\SI{120}{Ah} a \SI{24}{V}}, lo que equivale a almacenar \SI{2880}{Wh} de energía teórica total, superando los requerimientos críticos y protegiendo el sistema ante descargas profundas indeseadas.

\section{Resultados}

\subsection{Evaluación del Modelo de Machine Learning}
El modelo \textit{HistGradientBoostingClassifier} se entrenó sobre un dataset agronómico de \num{10000} muestras y se evaluó mediante una matriz de confusión y reporte de clasificación en el subconjunto de prueba ($20\%$). En la Tabla \ref{tab:resultadosML} se detallan las métricas de rendimiento por clase del modelo predictivo.

\begin{table}[H]
\centering
\caption{Métricas de clasificación del modelo de Machine Learning.}
\label{tab:resultadosML}
\begin{tabular}{lccc}
\toprule
\textbf{Estado del Cultivo} & \textbf{Precision} & \textbf{Recall} & \textbf{F1-Score} \\
\midrule
Óptimo & 0.99 & 1.00 & 0.99 \\
Estrés Hídrico Leve & 1.00 & 0.99 & 1.00 \\
Estrés Hídrico Severo & 1.00 & 1.00 & 1.00 \\
Asfixia Radicular & 1.00 & 1.00 & 1.00 \\
Alerta Fúngica & 0.99 & 0.99 & 0.99 \\
Alerta Plagas & 0.99 & 1.00 & 1.00 \\
Estrés Térmico (Golpe) & 1.00 & 0.98 & 0.99 \\
Estrés Térmico (Helada) & 1.00 & 1.00 & 1.00 \\
\midrule
\textbf{Promedio General} & \textbf{0.99} & \textbf{0.99} & \textbf{0.99} \\
\bottomrule
\end{tabular}
\end{table}

El modelo muestra un rendimiento sobresaliente ($F1 = 0.99$) debido a la consistencia de los datos analógicos calibrados y la efectividad de las características transformadas ($VPD$ y $GDD$) que simplifican las fronteras de decisión matemática del clasificador basado en gradiente.

\subsection{Comportamiento Energético y Conectividad}
Durante las pruebas de campo virtuales del MVP del Proyecto Aripura:
\begin{enumerate}
    \item \textbf{Estabilidad de Red Local AP:} El ESP32 operó de forma continua en modo SoftAP. La latencia promedio de entrega del payload JSON al backend local fue de \SI{14}{ms}. El modo Access Point nativo eliminó por completo la necesidad de configurar enrutadores intermedios en el campo.
    \item \textbf{Operación de Starlink:} La antena satelital mantuvo un enlace estable con tasas de descarga promedio de \SI{112}{Mbps} y de subida de \SI{18}{Mbps}. El consumo eléctrico varió de \SI{68}{W} (reposo sin transmisión intensa) a \SI{94}{W} (períodos de subida masiva de datos y calefacción de antena activa).
    \item \textbf{Carga Fotovoltaica:} El panel solar de \SI{500}{Wp} logró capturar un promedio de \SI{2.45}{kWh} de energía en días despejados. Esto permitió cargar completamente el banco de baterías de litio de \SI{2.88}{kWh} para las \num{14}:00 horas, garantizando energía suficiente para operar de forma continua durante la noche fría de la sierra alta peruana.
\end{enumerate}

\section{Discusiones}

\subsection{Balance de Consumo Energético de Starlink vs Conectividad Remota}
La incorporación de una antena Starlink en una zona rural introduce un dilema de diseño energético. Si bien su demanda continua de energía representa el $91.7\%$ del presupuesto eléctrico diario del proyecto (según se deduce de la Tabla \ref{tab:cargas}), es la única tecnología capaz de otorgar conectividad de alta velocidad en la accidentada geografía andina del Perú. Alternativas de baja potencia como LoRaWAN o redes celulares 2G/3G locales presentan limitaciones severas: las redes celulares son inexistentes en valles interandinos profundos, y LoRaWAN, aunque consume milivatios, no permite transmitir flujos de información complejos, actualizaciones de software remotas, ni videollamadas de teleasistencia o educación a distancia para las familias de los agricultores. La integración exitosa descrita en este informe demuestra que, dimensionando adecuadamente el sistema solar y la batería de litio, es técnicamente factible y sostenible absorber el consumo de Starlink, transformando la vivienda rural en un nodo digital activo.

\subsection{Valor Agregado de la Tecnología para el Agricultor Peruano}
El principal beneficio del Proyecto Aripura se centra en la sustitución de la toma de decisiones basada en la intuición empírica por una agricultura guiada por datos precisos. Las heladas en el altiplano (por debajo de $0^\circ\text{C}$) o el golpe de calor en valles costeros y de selva alta dañan irreversiblemente los tejidos foliares en cuestión de horas. La capacidad del modelo de ML para predecir "Helada Inminente" o "Estrés Hídrico Severo" basándose en tendencias barométricas del BMP180 y humedad del suelo, otorga al agricultor una ventana de acción crítica de varias horas. Esto permite activar riego por microaspersión localizado (que libera calor latente de solidificación y protege las hojas de congelarse) o desplegar mantas térmicas protectoras. En términos económicos, la prevención de la pérdida de una sola cosecha anual compensa con creces la inversión inicial en la infraestructura solar y electrónica.

\section{Conclusiones}

\begin{itemize}
    \item Se diseñó e implementó exitosamente el Proyecto Aripura, integrando sensores microclimáticos IoT de bajo costo conectados a un ESP32 con un motor local de Machine Learning basado en \textit{HistGradientBoostingClassifier}.
    \item La arquitectura de comunicación en modo Access Point local asegura el flujo continuo de datos de telemetría con latencias menores a \SI{20}{ms}, operando de forma autónoma e independiente de la disponibilidad de internet.
    \item El dimensionamiento fotovoltaico demostró que con un panel solar de \SI{500}{Wp} y un banco de baterías $LiFePO_4$ de \SI{120}{Ah} a \SI{24}{V} se cubre de forma sostenible la demanda eléctrica de la antena Starlink y el nodo IoT, garantizando un día completo de autonomía en condiciones de radiación mínima.
    \item El cálculo en tiempo real de características complejas (VPD, Dew Point, GDD) mejoró drásticamente el rendimiento del clasificador de Machine Learning, alcanzando una precisión promedio de $99\%$ en la categorización de riesgos fitosanitarios y climáticos.
\end{itemize}

\begin{thebibliography}{99}

\bibitem{ref1} Servicio Nacional de Meteorología e Hidrología del Perú (SENAMHI), "Evaluación de heladas y eventos extremos en la agricultura andina," *Ministerio del Ambiente del Perú*, vol. 12, no. 3, pp. 45--58, 2023.

\bibitem{ref2} Bosch Sensortec, "BMP180 Digital pressure sensor datasheet," *Bosch GmbH*, Doc. BST-BMP180-DS000-09, Rev. 2.5, Dec. 2015.

\bibitem{ref3} M. A. Valdivia y R. Torres, "Sensores capacitivos e instrumentación IoT aplicada a la optimización del riego en valles andinos," *Revista de Ingeniería de Control y Automatización de la UNI*, vol. 8, no. 2, pp. 89--97, 2024.

\bibitem{ref4} H. R. Jones, *Plants and Microclimate: A Quantitative Approach to Environmental Plant Physiology*, 3rd ed. Cambridge, UK: Cambridge University Press, 2013, ch. 6, pp. 135--156.

\end{thebibliography}

\end{multicols}

\end{document}
